\section{Cross-track diagnostics and ablation evidence}\label{sec:extended}
\subsection{Retrospective GUIDE reference}
The historical GUIDE program normalized evidence rows, excluded ambiguous alert membership from its local statistical grouping scope, formed bounded entity/time candidates, learned pair scores, capped connected components at 100 alerts, and ranked resulting groups. Its frozen v2 grouping plus v7 ranking reference was evaluated over 499 organization queues, 400,039 unambiguous alerts and 9,980 expert top-20 labels. Reported end-to-end Recall@20 was 49.82\%, P@10 61.58\%, nDCG@20 0.4371, group recoverability 69.83\%, and grouped-pair precision $269{,}832/298{,}438=90.4148\%$. It also produced 28,606 false grouped pairs and 4,841 mixed predicted groups. These retrospective data had previously exposed test labels, and source-field timing/authorization remained unresolved. The result is exploratory, cannot be called a fresh test, and is not an evaluation of the synthetic three-root verifier.

The GUIDE pair-recall denominator is all $159{,}186{,}107$ positive alert pairs in its defined eligible scope: $269{,}832/159{,}186{,}107\approx0.1695\%$. This very low exhaustive-pair recall can coexist with useful incident-level retrieval because the metrics have different units and large incidents dominate possible pair counts. Candidate positive-pair coverage was $699{,}790/159{,}186{,}107\approx0.4396\%$; an oracle connected 8,638 of 9,980 expert labels, whereas the actual fixed-group recovery was 69.83\%. The differences locate losses in candidate generation, scoring/merging and ranking without treating oracle connectivity as attained performance. Sixty-two expert incidents with more than 125 in-scope alerts cannot attain 80\% completeness under a hard 100-alert group cap; they remain in the denominator.

\begin{table*}[t]\centering\footnotesize\caption{Separate quantitative tracks and their denominators. Numbers across rows are not a common benchmark.}\label{tab:cross_track}
\begin{tabularx}{\textwidth}{@{}p{.22\textwidth}p{.18\textwidth}Y Y@{}}\toprule
Study & Scope & Headline R@20 & Safety/evidence boundary\\\midrule
Historical GUIDE v2+v7 & 499 queues, 9,980 labels & 49.82\% retrospective exploratory & 90.4148\% grouped-pair precision; 28,606 false pairs; no fresh label isolation\\
Controlled development & 10 queues, 200 opportunities & 158/200 verified base; 169/200 fallback & Synthetic only; fallback missed 170/200 development gate\\
Frozen matched holdout & 15 queues, 300 opportunities & 250/300 base; 260/300 fallback & Zero observed false pairs/mixed groups, simulator truth\\
External SOC & No compatible evaluated cohort & Not measured & No real-data prediction or production assessment\\\bottomrule
\end{tabularx}\end{table*}

\subsection{Rejected changes and failure mechanisms}
The controlled development sequence matters because it documents alternatives that looked plausible but did not pass safety or recall gates. A richer single-lineage claim variant yielded 2,580 false pairs and 22 mixed groups on its development cohort; independent verification removed observed false grouping but left Recall@20 at 158/200. The narrow endpoint substitution raised development recall to 169/200 and recoverability to 177/200, with 14 newly accepted genuine unions, yet missed its 170/200 selection gate. The unchanged candidate was later tested once on a fresh generator-family holdout under a documented protocol adjustment. This history weakens a strict preregistration claim and should accompany the positive result.

On the retrospective GUIDE track, 38 small-fragment heuristic merges included 25 false incident joins and added 168 false grouped pairs; candidate expansion added 49,455 possible pairs but no net complete/pure incident recovery. A later rank-boundary candidate reduced grouped-pair precision dramatically to 54.98\% when it changed pair grouping. A consensus split improved pair precision to 92.07\% but reduced Recall@20 to 45.13\%. A validation-selected fixed-group rank fusion reached 50.77\% R@20, 63.88\% P@10 and 0.4697 nDCG@20 on the already inspected 499-queue set; it is not a confirmatory result and does not alter the v2+v7 reference used here. These failed controls make the tradeoff between candidate coverage, purity and retrieval concrete.

Separate development-only synthetic studies changed ranking features, near-boundary ordering, evidence availability and hard-negative training. Table~\ref{tab:development} reports matched rows without implying a temporal improvement curve. In the hard-negative study, the learned ranker moved matched R@20 from 263/300 to 262/300 while the partition stayed fixed; pairwise ordering loss and official top-20 incident recall are different objectives. Under delayed, missing, derived-copy, common-root and contradictory telemetry, an unchanged verifier generally abstained more or recovered fewer incidents. No new independent telemetry source was available to justify a broader reconstruction rule. These observations support conservative verification but do not bound real-world false-merge probability.

\subsection{Interpretation of the holdout error matrix}
The token matrix in Table~\ref{tab:confusion} conditions on audited visible relationship-token decisions. Its 95.31\% genuine-token acceptance recall measures a different object than 86.67\% incident Recall@20. A single accepted token may change many within-group pairs and may or may not make a complete/pure incident; a true rejected token need not cause an incident miss. The paired incident table records 15 newly recalled and five lost recalls. The latter demonstrates that reconstruction can perturb ranking sufficiently to push a previously retrieved incident below rank 20 despite a net gain. The heatmap shows two queues with a one-incident retrieval loss and two P@10 declines. A zero false-positive cell on this small synthetic sample is a measured count, not a proof of universal specificity or attack resistance.
